% !TeX root = ../数模通用模板.tex
\section{问题分析}
四问均需要在供需平衡、储电范围和功率限制下确定购电安排。其递进关系体现为供需从确定量变为预测量、光伏预报从每日一次变为日内更新、电价从固定曲线变为实时波动。本文先建立能量与费用的公共表达，再分别处理各问的信息条件和运行要求。实际执行以已锁定的购电量和当槽实测供需为依据，连续更新储电量。
\subsection*{问题1分析}
给定典型日曲线后，问题一可以直接建立混合整数线性规划。首先在充放电互斥、日初日末储电量相同的基础上，增加净功率变化和模式最短运行时间设置，求最低购电费用；再允许费用在第一层结果基础上增加0.1\%，减少功率波动和启动次数。两个目标分层处理，使平稳运行的费用代价具有明确参照，最后由第二层同一条轨迹生成购电、充放电和指定时段结果。
\subsection*{问题2分析}
计划剩余与供电不足的费用不对称：未利用的计划电量仍需付费，缺口则产生5倍电价的紧急支出。负载和光伏分别采用直方图梯度提升树与极端随机树的原始预测等权融合，再进行近28日岭校准和前日负载偏差修正。保留整日误差的时间相关性，先用代表路径确定各槽共同充放电模式，再固定模式修正购电量，最后按真实供需执行。每日计划模式变化预算用于限制频繁改变运行方向，实际储电量和计划末模式跨日传递。
\subsection*{问题3分析}
新的光伏预报能改变对剩余供需的判断，但调整需要支付相应费用。午夜负载采用单一梯度提升树预测，光伏从0时起使用当时发布的官方预报；在6、12、18时重新校准光伏节点并修正负载偏差。利用历史同发布时间误差路径建立情景线性规划，各情景共享购电量，储能响应分别计算。仅对当日未执行区间重新规划，以实际储电量为起点；用下一次更新前后的不同短缺权重近似后续调整机会。最终调整量相对午夜原计划结算一次，下调部分退回原费用的一半。
\subsection*{问题4分析}
未来真实价格不提前进入决策，先用近28日已观察价格及周期特征进行加权岭回归。无日内预报的方案在价格模型中加入已发布的负载与光伏预测，并采用小时共同模式及固定模式购电修正；有日内预报的方案保留八维价格回归，在每次发布时更新价格并求解剩余时段的情景线性规划。两者均按实际价格结算。比较时报告完整方案的费用差，并说明预测器、信息权限与控制参数的差异，不将跨问差额解释为某一个因素的独立效果。
